
\section{One bridge compiles the normal-host reduction}\label{sec:compiler}
We extend the label refinement and rational-row tools of
\cite[Lemma 4.2 and Theorem B]{Entropy} from spatial strategies to normal
commuting strategies. The finite predicates are retained. The proofs
below specify where separation replaces Schmidt-range cancellation.

\subsection{The calibrated normal bridge}
Here FMS means \emph{flagged magic square}; Fu--Miller--Slofstra
are cited by name. A question or projection is \emph{synchronized}
when a binary event is perfectly matched with its opposite-party
copy; on faithful support it is a centralizer projection with its
standard commutant mirror. Here is the finite predicate, in zero-based indices.
Alice asks $R_0,R_1,R_2,T_0,T_1$ and Bob asks $C_0,C_1,C_2,U_0,U_1$.
A row answer is a flag $f$ or a triple of product $+1$.
A column answer is $f$ or a triple of product $+1,+1,-1$ in the
respective columns. At $(R_i,C_j)$ flags must match; nonflags agree
on their intersecting cell. Choose $\gamma_0=\rho_0=(+++)$ and
$\gamma_1=\rho_1=(+--)$. At $T_x$ replace $f,\gamma_x$ by signs $+,-$ and retain
the other three column-zero triples; $U_y$ does the same with
$f,\rho_y$ in row zero. At $(T_x,C_0)$ signs may match only $f,\gamma_x$,
and retained triples must equal Bob's triple. Impose the symmetric
rule at $(R_0,U_y)$. At $(T_x,U_y)$ two signs $a,b$ must obey
$ab=(-1)^{xy}$; pairs with a retained triple are accepted.
Unspecified pairs are accepted.

\begin{center}\small
\begin{tabular}{@{}p{.24\linewidth}p{.66\linewidth}@{}}
\toprule Question pair & Winning rule\\\midrule
$(R_i,C_j)$ & flags match; nonflags agree at cell $(i,j)$\\
$(T_x,C_0)$ & a sign pairs only with $f,\gamma_x$; retained triples match\\
$(R_0,U_y)$ & symmetric calibration at $\rho_y$\\
$(T_x,U_y)$ & two signs obey $ab=(-1)^{xy}$; a retained triple is accepted\\
Conditional copy & both inactive or both active; activity, flags and selected events match their binary names\\
Other pairs & all valid answer pairs accepted\\\bottomrule
\end{tabular}
\end{center}

Conditional copies also have an inactive answer; reject exactly one
inactive answer, and bind activity, flags and the stated magic events
to named synchronized binary questions.

The Pauli square of Mermin--Peres \cite{Mermin,Peres} is
\[
 \begin{pmatrix}XI&IX&XX\\ IZ&ZI&ZZ\\ XZ&ZX&YY\end{pmatrix}.
\]
Its four column-zero outcome projections are
$m_{00},m_{01},m_{10},m_{11}$ in the order
$(+++),(+--),(-+-),(--+)$.
On Alice, read $m_{00}$ from the retained answer of $T_1$ and the
other three from $T_0$; on Bob read them from $C_0$.
Bind each chosen occurrence to the opposite binary name.
Also bind positive $ZI$ from Alice's $R_1$ component 1 and Bob's
$C_1$ component 1, and positive $IX$ from Alice's $R_0$ component 1
and Bob's $C_1$ component 0. If $M_{\rm mag}=\sum m_{ab}$, put
\[
 x=I-2M_{\rm mag}+2E_{ZI,+},\qquad
 z=I-2M_{\rm mag}+2E_{IX,+}.
\]
These are centralizer reflections, identity off the magic support.
They transport $m_{00}$ to $m_{ab}$ by $x^az^b$.

For clarity we recall the arbitrary-normal bridge argument of
\cite[Lemma 2.2]{ReadType}. Matching makes the cells sharp centralizer
effects. Write the upper-left four cell reflections as $a,b,c,d$.
The row/column relations give
$[a,b]=[a,c]=[b,d]=[c,d]=0$,
$(ac)(bd)=-(ab)(cd)$, and $[ac,bd]=0$.
Canceling the outside $a,d$ gives $cb=-bc$; substitution into the
last commutator gives $ad=-da$.
The two commuting Pauli pairs therefore generate centralizer $M_4$.
In $\mathbb C^2\otimes\mathbb C^2$ the calibrated unit vectors are
$\mathfrak c_x=|+\rangle\otimes|x\rangle$ and
$\mathfrak r_y=|+\rangle\otimes|(-1)^y\rangle$, where
$|+\rangle=(|0\rangle+|1\rangle)/\sqrt2$ and
$|(-1)^y\rangle=(|0\rangle+(-1)^y|1\rangle)/\sqrt2$.
Choose a minimal magic projection $q$ and centralizer partial
isometries $u_x,v_y$ with initial projection $q$ and ranges the
projections onto $\mathfrak c_x,\mathfrak r_y$ (tensored with the
multiplicity identity). Choose their phases so that their overlap is
$u_x^*v_y=(-1)^{xy}q/\sqrt2$.
If $\Delta_x^A,\Delta_y^B$ are the original sign differences,
their pull-backs on the standard corner
$(f+q)\mathcal J(f+q)\mathcal J L^2(M)$ are
\[
 X_x=(f+u_x^*)\Delta_x^A(f+u_x),\qquad
 Y_y=\mathcal J(f+v_y^*)\mathcal J\Delta_y^B \mathcal J(f+v_y)\mathcal J.
\]
They are commuting self-adjoint contractions.
Their joint-support vector is
\[
 \theta_{xy}=(f+(-1)^{xy}q/\sqrt2)\xi_\varphi .
\]
Forbidden sign events give
$X_x\theta_{xy}=(-1)^{xy}Y_y\theta_{xy}$ and
$X_xY_y\theta_{xy}=(-1)^{xy}\theta_{xy}$.
The multiplier $f\pm q/\sqrt2$ is bounded invertible in the corner
centralizer, so $\theta_{xy}$ separates both corner algebras.
Commuting the opposite contraction through the first identity gives
$X_x^2\theta_{xy}=\theta_{xy}$ and $Y_y^2\theta_{xy}=\theta_{xy}$;
separation makes them reflections.
Put $\kappa=f+q/\sqrt2$. The three positive cells give
$X\kappa\xi_\varphi=Y\kappa\xi_\varphi$ and identify
$X_0=X_1=X$ and $Y_0=Y_1=Y$
by separation on $\kappa\xi$. The fourth vector is $Z\kappa\xi$, $Z=f-q$.
Its equation is $XZ\kappa\xi=-YZ\kappa\xi=-ZY\kappa\xi=-ZX\kappa\xi$,
since opposite algebras commute. Separation gives $XZ=-ZX$.
Thus $X=v+v^*$, $v=qXf$, $v^*v=f$, $vv^*=q$.
The mirror equation means that $X$ centralizes
$\omega(\cdot)=\varphi((f+q/2)\cdot)$.
The bounded centralizer perturbation formula of Pedersen--Takesaki \cite[Lemma 4.5]{PT}
$\sigma_t^\omega(y)=k^{it}\sigma_t^\varphi(y)k^{-it}$,
$k=f+q/2$, gives on the $q$--$f$ block
\begin{equation}\label{eq:bridge}
 \sigma_t^\varphi(v)=2^{it}v,\qquad \varphi(q)=2\varphi(f).
\end{equation}
The last equality is KMS. Multiplication by the magic matrix units
gives all four transports. If one sector vanished, the off-diagonal
reflection would kill the other; hence on a nonzero active corner
both sectors are nonzero. No finite multiplicity is used.

\begin{lemma}[Commuting label refinement]\label{lem:refine}
Retain the old FMS questions. For a synchronized projection $Z_0$, add
$T_0[Z_0]$ with answers $(a,z_0)$, use the old rule after forgetting
$z_0$ against unchanged $U_0,U_1,C_0$, and match $z_0$ to Bob's
binary $Z_0$. Add the symmetric refined $U_0[Z_0]$ against unchanged
$T_0,T_1,R_0$. Retain activity matching.
In every perfect normal commuting strategy the coarse effects equal
the original effects and commute with $Z_0$. Conversely, commuting
labels give perfect refinements. Separate labels need not commute.
\end{lemma}
\begin{proof}
Calibration fixes the retained effects and the sign support.
Pull the new coarse sign difference $X'$ to the same corner.
Against the unchanged opposite sign it satisfies
$X'\theta_{0y}=Y_y\theta_{0y}=X_0\theta_{0y}$.
The separating vector above gives $X'=X_0$.
Thus its coarse effects $E_a$ are the old projections.
For the fine effects $Q_{a,z_0}$, label matching gives
$\sum_aQ_{a,1}=Z_0$ and $\sum_aQ_{a,0}=I-Z_0$ on faithful support.
Positivity gives $0\leq Q_{a,1}\leq E_a,Z_0$ and
$0\leq Q_{a,0}\leq E_a,I-Z_0$.
Hence $E_aZ_0=Q_{a,1}=Z_0E_a$, using completeness of each marginal.
This proves commutation, and the symmetric argument handles Bob.
Conversely use the joint products $E_aZ_0,E_a(I-Z_0)$.
Each label has its own refined question with the same unchanged
opposite signs, so there is no cross-label commutation constraint.
\end{proof}

\subsection{Graphs turn products into rational rows}
All projections here are synchronized, hence centralizer projections.
For orthogonal $F_i,F_j$, $F_{ij}=F_i+F_j$, impose
\begin{equation}\label{eq:graph}
 G+K_G=F_{ij},\quad H_G+K_H=F_{ij},\quad 5H_G+F_{ij}=3F_i+4G .
\end{equation}
On $F_{ij}$, $Z=F_i-F_j$ and $X=2G-F_{ij}$ are reflections.
Since $2H_G-F_{ij}=(3Z+4X)/5$, squaring gives $ZX+XZ=0$.
Therefore
\[
 X=\begin{pmatrix}0&u_{ij}\\u_{ij}^*&0\end{pmatrix},
 \qquad u_{ij}^*u_{ij}=F_j,\quad u_{ij}u_{ij}^*=F_i.
\]
For a triple partition $F_1+F_2+F_3=I$, impose these graphs on
all pairs and add a projection $\Pi_\triangle$ with
$3\Pi_\triangle=2(G_{12}+G_{23}+G_{13})-I$.
Its diagonal blocks are $I/3$, off-diagonal blocks $u_{ij}/3$.
The $(1,3)$ block of $\Pi_\triangle^2=\Pi_\triangle$ gives $u_{13}=u_{12}u_{23}$.
Conversely this identity makes the entire block matrix a projection.
A reference identity triangle gives the scalar $M_3$ frame.
We call this the \emph{qutrit frame}: its three diagonal projections
and identity graph projections specify the matrix units of a common
scalar $M_3$, so all application operators act on its multiplicity
space.
Triangles $(U,1,U)$ and $(1,U,U)$ identify all placements of a unitary;
a triangle $(U,V,W)$ then certifies $W=UV$.
An inverse is a product equal to the identity. Take proper prefixes
of the finite input word list to make all required products explicit.

To bind a prescribed synchronized projection $E$, introduce
projections $EF_i,NF_i$ with
$EF_i+NF_i=F_i$, $\sum_iEF_i=E$, $\sum_iNF_i=I-E$.
For each $ij$ put $F_{ij}=F_i+F_j$ and add a joint partition
$B^{ij}_{00}+B^{ij}_{01}+B^{ij}_{10}+B^{ij}_{11}=F_{ij}$ with marginals
$EF_i+EF_j$ and $G_{ij}(1)$; the fifth, inactive outcome is $I-F_{ij}$.
Define $G_{ij}(s_E)=B^{ij}_{11}+B^{ij}_{00}$, the equal-bit event inside $F_{ij}$.
Sharp marginals make these actual joint products.
They force $E$ to commute with the frame and give
$s_E=2e-I$ when $E=I_3\otimes e$.
Thus the parameter is the specified old event.

For a real rational polynomial $p=\sum_wq_wU_w$, the row
\begin{equation}\label{eq:offdiag}
 2\sum_wq_wG_{12}(U_w)
 -\left(\sum_wq_w\right)(F_1+F_2)=0
\end{equation}
has off-diagonal blocks $p,p^*$, so enforces even a non-self-adjoint
equation $p=0$. These arguments are bounded operator identities
in the centralizer, which is finite: $\varphi|_{M_\varphi}$ is a
faithful normal trace.

\begin{lemma}[Rational rows on normal commuting hosts]\label{lem:rows}
Given finitely many synchronized projections and real rational linear
identities among them, a finite Boolean extension enforces those
identities on faithful normal commuting supports. Every realization
extends after finite matrix amplification in standard form, preserving
the old nonsynchronous tests and the host type.
\end{lemma}
\begin{proof}
We give the finite extension and both directions.
Certify scalar $M_{d_{\max}}$ by synchronized magic squares, with
$d_{\max}$ a power of four. Joint questions between axes of distinct
squares force their commutation; the Pauli calculation above identifies
the tensor matrix algebra. Joint questions between every old projection
and every axis force the old projections into its commutant.
Thus old projections are $I\otimes e_i$ and scalar matrix units are
centralizer elements. Bind global $X$ and $Z$ bit measurements to the axes.

For a row $\mathsf S_{\rm row}=\alpha_0I+\sum_i\alpha_i e_i$, clear denominators with $\ell_{\rm row}$.
Using $s_i=2e_i-I$, write $2\ell_{\rm row}\mathsf S_{\rm row}$ as a list of signed reflections
$V_j\in\{\pm I,\pm s_i\}$ with sum $2\ell_{\rm row}\mathsf S_{\rm row}$ and first slot $V_0=I$.
Add cancelling $I,-I$ pairs as needed and pad to a power of four $d$.
Its length can be even because the coefficient sum
$2\ell_{\rm row}\alpha_0+2\sum_i\ell_{\rm row}\alpha_i$ is even.
For each $j\ne0$, joint scalar parity questions certify
\[
 \Pi_j^\pm=\tfrac12(|0\rangle\pm|j\rangle)
                 (\langle0|\pm\langle j|).
\]
Explicitly use $X^j$ and a basis of the commuting $Z$ parities
$j^\perp$; select zero syndrome and the indicated sign.
Joint questions of $e_i$ with this three-outcome scalar measurement
certify the graph coarse events
$G_j^\pm=\frac12\left(\begin{smallmatrix}I&\pm V_j\\
\pm V_j^*&I\end{smallmatrix}\right)$ on coordinates $0,j$.
Here $V_j$ is self-adjoint. The displayed $2\times2$ matrix acts
on $\operatorname{span}\{|0\rangle,|j\rangle\}$ tensored with the
old space and is zero on the remaining scalar coordinates.

Add a synchronized partition $(\Pi_t^{\rm row})_{t\in\mathbb F_2^{\log_2d}}$.
Forbid $\Pi_t^{\rm row}$ with $G_j^-$ when
$\chi_t(j)=(-1)^{t\cdot j}=1$, and with $G_j^+$ otherwise.
A forbidden pair with mirror $\mathcal JG\mathcal J$ gives
$\Pi_t^{\rm row}G\xi_\varphi=0$, since $\mathcal JG\mathcal J\xi_\varphi=G\xi_\varphi$;
separation gives $\Pi_t^{\rm row}G=0$. Thus a vector in $\ran \Pi_t^{\rm row}$ has
coordinates $v_j=\chi_t(j)V_jv_0$.
Define isometries
\[
 I_tu=d^{-1/2}(\chi_t(j)V_ju)_j,\qquad P_t=I_tI_t^*.
\]
Character orthogonality gives $I_t^*I_u=\delta_{tu}I$ and
$\sum_tP_t=I$, without commutation of distinct $V_j$.
Hence $\Pi_t^{\rm row}\leq P_t$; the two partitions sum to $I$, so $\Pi_t^{\rm row}=P_t$.
Finally forbid $\Pi_0^{\rm row}$ with the scalar projection
$\Pi_d=|+_d\rangle\langle+_d|$, the all-positive global $X$ event.
Then
\[
 I_0^*\Pi_dI_0=d^{-2}\Big(\sum_jV_j\Big)^*
                         \Big(\sum_jV_j\Big)
\]
vanishes exactly when $\mathsf S_{\rm row}=0$.

Conversely amplify by $M_{d_{\max}}$ with normalized trace.
Use the actual scalar Paulis, joint products and projections $P_t$.
All are centralizer projections. Bob uses $\mathcal J_{\rm new}X\mathcal J_{\rm new}$
for each projection $X$. For a centralizer word $xy$,
$\mathcal J(xy)\mathcal J\xi_\varphi=(xy)^*\xi_\varphi$; adjunction supplies the
reversed order. Thus Bob's products are the actual transferred
products, even when the old projections do not commute.
Every structural forbidden event is zero, and the final event is
zero precisely for a zero row. Old measurements act trivially on
the new matrix factor, preserving all old probabilities.
Finite matrix amplification preserves all indicated factor types.
\end{proof}

\subsection{The complete one-bridge assembly}
Use five orthogonal synchronized corners
$e_0,m_{00},m_{01},m_{10},m_{11}$ covering $I$, with one FMS copy
whose flag is $e_0$ and magic support is their other four corners.
Put $e_1=m_{00}$. The actual bridge $v:e_0\to e_1$ obeys
\eqref{eq:bridge}. KMS and the cover give $\varphi(e_0)=1/9$ and
$\varphi(m_{ab})=2/9$. These masses are independent of the root model.

\begin{figure}[ht]
\centering
\[
 \underbrace{e_0}_{1/9}
 \ \xrightarrow[\sigma_t^\varphi(v)=2^{it}v]{v}\
 \boxed{\underbrace{m_{00}=e_1}_{2/9}\quad
 m_{01}\quad m_{10}\quad m_{11}}
 \qquad\text{each magic corner has mass }2/9.
\]
\caption{One calibrated bridge. All copies of transported labels use
this same $v$; the box is the four-dimensional magic support.}
\end{figure}

Bind the five corners and the two positive $ZI,IX$ parameters to the
qutrit frame. Introduce the two axis unitaries $x,z$ with the explicit
formulas above. A \emph{typed label} is an operator name together with its required
class (projection, positive contraction or unitary).
The transported typed label list is exactly
\begin{equation}\label{eq:labels}
 Q,\ C,\ T,\ L,\ \{U_g:g\text{ is a generator of }\mathcal G_e\}.
\end{equation}
Encode a projection by a reflection, a positive contraction by
$u_C=2C-I+2i\sqrt{C(I-C)}$, and a unitary directly.
A \emph{unitary name} is a formal generator for the chosen unitary
encoding, including separate names for separate summands when needed.
For each resulting unitary name $k$ introduce copies $u_{r,k}$,
$r=0,1$, with $(u_{r,k}-I)(I-e_r)=0$.
The \emph{decoded copies} are the actual corner operators recovered
from these encodings by the following formulas, respectively:
$e_r(I+u_{r,Q})/2$,
$e_r(2I+u_{r,C}+u_{r,C}^*)/4$, and $e_ru_{r,k}$.
Require $u_{r,Q}^2=I$.
Introduce a block unitary
\[
 V_k=e_0u_{0,k}+
 \sum_{a,b}x^az^b e_1u_{1,k}z^{-b}x^{-a}.
\]
Refine the retained FMS signs separately for the six frame projections
\[
 F_1,\quad F_2,\quad F_3,\quad G_{12}(1),\quad G_{13}(1),\quad G_{23}(1),
\]
and for every $G_{12}(V_k)$.
Lemma~\ref{lem:refine} first makes $v$ preserve the
frame, so $v=I_3\otimes\widetilde v$.
The graph-label commutators in the target--source block give
\begin{equation}\label{eq:transport}
 vu_{0,k}=u_{1,k}v.
\end{equation}
Thus the same actual calibrated transport works for every label.
Transport of decoded adjoints and polynomials follows too.
$P$ is root-only, encoded by $P_0=e_0(I+u_P)/2$ with
$u_P^2=I$ and $(u_P-I)(I-e_0)=0$.

Introduce two centralizer contractions $b_i=e_1b_ie_0$.
For compilation encode each as an average of two global unitaries.
Extract $B_i=v^*b_i$ on $e_0Me_0$.
Their modular eigenvalue is $1/2$ because that of $b_i$ is one.
The complete endpoint translation is
\begin{center}\small
\begin{tabular}{@{}ll@{}}
\toprule Root relation & Centralizer polynomial\\\midrule
$B_i^*B_j=\delta_{ij}(I-P)$
& $b_i^*b_j=\delta_{ij}(e_0-P_0)$\\
$\sum_iB_iB_i^*=I-Q$ & $\sum_i b_ib_i^*=e_1-Q_1$\\
$CB_i=\frac12B_iC$ & $C_1b_i=\frac12b_iC_0$\\
$U_gB_i=B_iU_g$ & $(U_g)_1b_i=b_i(U_g)_0$\\
$TB_i=B_iT$ & $T_1b_i=b_iT_0$\\
$LB_i=B_i aLa^*$ & $L_1b_i=b_i a^{[0]}L_0(a^{[0]})^*$\\
\bottomrule
\end{tabular}
\end{center}
Every remaining relation is a root centralizer polynomial:
the finite relators of $\mathcal G_e$, $C_0Q_0=Q_0$, $[C_0,P_0]=0$,
$[U_g,C_0]=[T_0,C_0]=0$,
$L_0Q_0=c^{[0]}Q_0$, $[L_0,C_0]=[L_0,P_0]=0$,
and the terminal equation.
Expand its root projection as $P=(I+s_P)/2$.
Here $a^{[0]},c^{[0]}$ mean the source words for the embedded $a,c$. If $W=[b,L]$, its non-self-adjoint polynomial is explicitly
\[
 p=\tfrac12\bigl(TWT^*+W+s_PTWT^*+s_PW\bigr).
\]
Use \eqref{eq:offdiag} for these words and their prefixes.
In global scaffold coordinates multiply by $e_0$ and expand
$e_0=(I+s_{e_0})/2$ as well. This forces $p=0$ and $p^*=0$ via
the two off-diagonal blocks; it never assumes $PW$ self-adjoint.
Expand all other typed polynomials in the same manner, apply
graphs and triangles, then Lemma~\ref{lem:rows}. Retain all FMS and
refinement questions. Sample all padded question pairs uniformly,
reject the listed events, and accept unspecified pairs.
This is an effective finite rational Boolean game, denoted $G_e$.

\subsection{Soundness and simultaneous completeness}
Soundness uses the nonzero root corner with normalized state
$9\varphi|_{e_0Me_0}$. Equations \eqref{eq:transport} and the table
give all root relations as bounded operator identities.
For example $C_0v^*=v^*C_1$ gives
$C_0B_i=v^*C_1b_i=B_iC_0/2$.
Conjugating the range equation by $v^*$ uses $v^*Q_1v=Q_0$.
The corner of a semifinite host is semifinite, so
Lemmas~\ref{lem:selector} and \ref{lem:reader} prove soundness.

Conversely take any faithful normal root realization $(M,\varphi)$.
On $M_5\vN M$ use the state
\begin{equation}\label{eq:scaffold}
 \varphi_{\rm sc}=\operatorname{diag}(1,2,2,2,2)/9\ \otimes\varphi.
\end{equation}
Use the identity corner maps $v=e_{10}\otimes I$, the ideal real
FMS effects, and $b_i=e_{10}\otimes B_i$.
The product modular group makes each $b_i$ centralizer because
the two modular eigenvalues $2$ and $1/2$ cancel.
The whole centralizer is finite, with its faithful trace
$\varphi_{\rm sc}$.
For any contraction $x$ in a finite von Neumann algebra, the polar
partial isometry extends to a unitary $u$: equivalent initial and
final projections have equivalent complements in a finite algebra.
Then
\[
 u_\pm=u(|x|\pm i\sqrt{I-|x|^2}),\qquad x=(u_++u_-)/2.
\]
This supplies the exact encodings of both $b_i$ at once.
Use $2Q-I,2P-I$ and the displayed functional-calculus encoding of
$C$, copying the same choices to each required corner.
Amplify by the scalar qutrit frame. All block labels commute with
the ideal FMS signs, so use their actual joint products for refinement.
Separate labels use separate questions.
Finally apply the standard-form construction in
Lemma~\ref{lem:rows}, with Bob's effects defined by its conjugation
$\mathcal J_{\rm new}$. Formula \eqref{eq:scaffold} and the ideal real FMS
effects give exactly the five-dimensional FMS probabilities, including
the nonsynchronous signs. All graph and endpoint identities hold
simultaneously. A finite-algebra root stays finite; after tensoring
with $R_{\rm hyp}$ a $\II_1$ host exists. A $\III_{1/2}$ root stays
$\III_{1/2}$ under the finite amplifications.

\Needspace{10\baselineskip}
\subsection{Chained bridges}
\begin{lemma}[Chained bridges]\label{lem:chain}
For any fixed $k\geq1$, the conditional FMS predicates, individual
label refinements and rational rows above compile the following
finite root specifications on arbitrary normal commuting supports.
The centralizer labels have prescribed types (projection, positive
contraction, contraction or unitary). Each edge $A_\alpha$ has modular
ratio $2^{-k}$, with isometry and range equations and endpoint
relations $L_{\rm end}A_\alpha=A_\beta B_{\rm end}$ or $L_{\rm end}A_\alpha=0$; $L_{\rm end},B_{\rm end}$ are real rational
polynomials in the typed centralizer labels (for example $B_{\rm end}=C/2$). Other relations are root
centralizer polynomials with real rational coefficients.
Every perfect strategy extracts that root specification on a
nonzero corner. Every faithful normal root realization extends
simultaneously to a perfect strategy after finite amplification,
preserving its factor type and, for a spatial finite-entropy model,
finiteness of entropy. Only labels occurring at a left endpoint (including a range-equation
label) need transport; transporting further labels is permitted and
changes neither direction.
\end{lemma}
\begin{proof}
Use the synchronized partition
\[
 e_0+\sum_{g=1}^k \mathsf M_g=I,\qquad
 \mathsf M_g=\sum_{\epsilon,\delta\in\{0,1\}}m_{g,\epsilon\delta},\qquad e_g=m_{g,00}.
\]
For bridge $g$ use a conditional FMS copy with flag $e_{g-1}$,
magic support $\mathsf M_g$, and activity $a_g=e_{g-1}+\mathsf M_g$.
The inactive effect is $I-a_g$. Match activity on every conditional
question and bind the flag and four chosen magic events to their
named binary projections, in both party orientations. The
flag/four-magic/inactive column partition then has precisely these
six marginals. Different copies have disjoint question names;
no agreement between their nonsynchronous signs is imposed.

The one-bridge proof applies in each $a_gMa_g$. Indeed $a_g$ is
centralizer, so its normalized restriction has the restricted
modular group. It gives $w_g^*w_g=e_{g-1}$, $w_gw_g^*=e_g$ and
$\sigma_t^\varphi(w_g)=2^{it}w_g$. These equations extend by zero
off $a_g$ and imply
\[
 v^{(k)}=w_k\cdots w_1,\quad
 (v^{(k)})^*v^{(k)}=e_0,\quad
 v^{(k)}(v^{(k)})^*=e_k,\quad
 \sigma_t^\varphi(v^{(k)})=2^{kit}v^{(k)}.
\]
This is an extracted operator, never a centralizer word submitted
to the graph compiler. KMS and the magic matrix units give
\[
 \varphi(m_{g,\epsilon\delta})=2^g\varphi(e_0),\qquad
 Z_k^{\rm sc}=1+4\sum_{g=1}^k2^g,\qquad \varphi(e_0)=(Z_k^{\rm sc})^{-1}.
\]
If an active corner vanishes its flag and all four magic corners
vanish together; otherwise the local bridge makes all five
nonzero. The recurrence therefore holds also before nonvanishing
is known, and the cover proves that $e_0\ne0$.

Bind every corner and every positive Pauli parameter to the same
qutrit frame. Let $\Lambda$ be the unitary encodings of the labels to be transported,
including every label used on a left endpoint: use a unitary directly,
a reflection for a projection, the displayed functional calculus for a
positive contraction, and two unitaries for a general contraction.
Introduce the axis labels $x_g=I-2\mathsf M_g+2E_{g,ZI,+}$ and
$z_g=I-2\mathsf M_g+2E_{g,IX,+}$ of each bridge as in
Section~\ref{sec:compiler}. For each $\ell_{\rm lab}\in\Lambda$
and $r=0,\ldots,k$, impose
$(u_{r,\ell_{\rm lab}}-I)(I-e_r)=0$. Introduce a unitary label $V_{g,\ell_{\rm lab}}$ and impose
\[
 V_{g,\ell_{\rm lab}}=I-a_g+e_{g-1}u_{g-1,\ell_{\rm lab}}
  +\sum_{\epsilon,\delta}x_g^\epsilon z_g^\delta e_gu_{g,\ell_{\rm lab}}z_g^{-\delta}x_g^{-\epsilon}.
\]
The orthogonal block supports make this a unitary. At bridge $g$
apply Lemma~\ref{lem:refine} separately to the six frame projections
and to each $G_{12}(V_{g,\ell_{\rm lab}})$. There is one refinement of each
retained sign question per label per bridge, on each party, always
against the unchanged opposite questions. There is no joint
refinement of distinct labels. The frame refinements give
$w_g=I_3\otimes\widetilde w_g$. In the selected magic--flag block
the graph commutator is exactly
\[
 w_gu_{g-1,\ell_{\rm lab}}=u_{g,\ell_{\rm lab}}w_g.
\]
Induction gives $(L_{\rm end})_kv^{(k)}=v^{(k)}(L_{\rm end})_0$ for every decoded transported
label and its adjoints and words. Labels used only at the source
need one source encoding and no refinement.

For each edge introduce a centralizer contraction
$b_\alpha=e_kb_\alpha e_0$ and extract
$A_\alpha=(v^{(k)})^*b_\alpha$. Compile the centralizer equations
\[
 b_\alpha^*b_\beta=\delta_{\alpha\beta}e_0,\qquad
 \sum_\alpha b_\alpha b_\alpha^*=e_k-(P_\circ)_k,
 \qquad (L_{\rm end})_kb_\alpha=b_\beta (B_{\rm end})_0
\]
for the corresponding root equations; the zero endpoint uses zero
on the right. The equations with a specified initial defect use
that defect in the first right side. More generally a range equation
$\sum_\alpha A_\alpha A_\alpha^*=R_{\rm ran}$, with
$R_{\rm ran}$ a transported centralizer polynomial, is compiled as
$\sum_\alpha b_\alpha b_\alpha^*=(R_{\rm ran})_k$.
For the Catalan input this is $R_{\rm ran}=I-BB^*-P$; its
centralizer contraction $B$ and projection $P$ are already in the
transported list. Multiplying by $v^{(k)}$ or
its adjoint proves both directions, since its initial and final
projections are $e_0,e_k$. Every root-only polynomial stays at
$e_0$. In particular the edge has modular ratio $2^{-k}$ and all
labels refer to the same extracted root operators. Graphs,
triangles and Lemma~\ref{lem:rows} enforce these finitely many
centralizer polynomials without changing the conditional tests.
This proves commuting soundness.

For completeness take a faithful normal root $(N,\psi)$ and on
$\mathbb M_{1+4k}\bar\otimes N$ put
\[
 \psi_{\rm sc}=(Z_k^{\rm sc})^{-1}\operatorname{diag}
 (1;\,2^{\oplus4};\,4^{\oplus4};\,\ldots;\,(2^k)^{\oplus4})
 \otimes\psi.
\]
Choose the actual matrix corner maps for $w_g$, the ideal real FMS
effects on $a_g$, and $I-a_g$ for every inactive answer.
Take $b_\alpha=|e_k\rangle\langle e_0|\otimes A_\alpha$.
Its modular ratio is $2^k2^{-k}=1$, so it lies in the finite
centralizer. Its polar partial isometry therefore has a unitary
completion, giving its two-unitary encoding. Choose the root
centralizer encodings once and copy them identically to the required
corners. All $V_{g,\ell_{\rm lab}}$ then commute with the FMS effects at their
own bridge, even though adjacent activities overlap at $e_g$.
Use their actual joint products for the separate refinements.
Append the common scalar frame and row factors and use standard
commutant mirrors as in Lemma~\ref{lem:rows}. All old conditional
probabilities and every added row are simultaneously correct.
Only finite matrix factors have been appended. Their state entropy
is finite; in the spatial case it adds to the root entropy.
This proves completeness and the stated preservation properties.
\end{proof}
